\documentclass{article}
\pdfinfoomitdate=1
\pdftrailerid{}
\usepackage{booktabs}
\usepackage{tabularx}
\usepackage{hyperref}
\usepackage{graphicx}
\usepackage{float}
\usepackage{array}
\usepackage{makecell}
\usepackage{ragged2e} % for RaggedRight in X columns
\usepackage{tabularray}
\usepackage{tabularx,booktabs,array,xcolor}
\usepackage{longtable}
\newcolumntype{L}[1]{>{\raggedright\arraybackslash}p{#1}}
\definecolor{mylight}{RGB}{255,255,180} % optional highlight

\newcolumntype{Y}{>{\RaggedRight\arraybackslash}X}
\newcolumntype{C}[1]{>{\centering\arraybackslash}p{#1}}

\hypersetup{
    colorlinks=true,       % false: boxed links; true: colored links
    linkcolor=red,          % color of internal links (change box color with linkbordercolor)
    citecolor=green,        % color of links to bibliography
    filecolor=magenta,      % color of file links
    urlcolor=cyan           % color of external links
}

\title{Hazard Analysis\\\progname}

\author{\authname}

\date{}

%% Comments

\usepackage{color}

% \newif\ifcomments\commentstrue %displays comments
\newif\ifcomments\commentsfalse %so that comments do not display

\ifcomments
  \newcommand{\authornote}[3]{\textcolor{#1}{[#3 ---#2]}}
  \newcommand{\todo}[1]{\textcolor{red}{[TODO: #1]}}
\else
  \newcommand{\authornote}[3]{}
  \newcommand{\todo}[1]{}
\fi

\newcommand{\wss}[1]{\authornote{magenta}{SS}{#1}}
\newcommand{\plt}[1]{\authornote{cyan}{TPLT}{#1}} %For explanation of the template
\newcommand{\an}[1]{\authornote{cyan}{Author}{#1}}

%% Common Parts

\newcommand{\progname}{RoCam} % PUT YOUR PROGRAM NAME HERE
\newcommand{\authname}{Team \#3, SpaceY
    \\ Zifan Si
    \\ Jianqing Liu
    \\ Mike Chen
    \\ Xiaotian Lou} % AUTHOR NAMES                  

\usepackage{hyperref}
\hypersetup{colorlinks=true, linkcolor=blue, citecolor=blue, filecolor=blue,
    urlcolor=blue, unicode=false}
\urlstyle{same}

\usepackage{indentfirst}
\usepackage{graphicx}

\usepackage{titling}

\pretitle{\begin{center}\includegraphics[width=0.5\textwidth]{../../assets/logo/black.png}\\[0.75em]\LARGE}
        \posttitle{\par\end{center}}

\usepackage[letterpaper, portrait, margin=1in]{geometry}

\usepackage{placeins}

\begin{document}

\maketitle
\thispagestyle{empty}

~\newpage

\begin{table}[hp]
  \caption{Revision History} \label{TblRevisionHistory}
  \begin{tabularx}{\textwidth}{llX}
    \toprule
    \textbf{Date} & \textbf{Developer(s)} & \textbf{Change} \\
    \midrule
    Date1         & Shike Chen            & Initial draft   \\
    Date2         & Shike Chen            & Detailed draft  \\
    \bottomrule
  \end{tabularx}
\end{table}

~\newpage

\tableofcontents

~\newpage

\wss{You are free to modify this template.}

\section{Introduction}
Hazard is generally defined as a harm or potential harm or effect caused by a
system which may lead to personal injury, property damage, environmental damage
etc. As an integrated system with both hardware and software components,
hazards may arise due to improper design, malfunction of the system, unexpected
user behavior, etc. Therefore, it is critical to identify potential hazards and
mitigate them in the early stage of the project. In the case of the RoCam
project, the potential hazards could cause physical injury of the user or
malfunction of a larger system.

\section{Scope and Purpose of Hazard Analysis}

The hazard analysis will identify any direct or potential hazards caused by the
system which includes hazardous materials, design, user safety, etc. The hazard
analysis will also address the general hazardous effects of this system which
includes hazards effects in the events of malfunction, misuse or environmental
damage. However, since the project is open source, any use case that is beyond
the scope of this project will not be included in this hazard analysis.

Such hazard analysis will not only used to identify potential hazards caused by
the system, but also to minimize risks related in unforeseen cases. The purpose
of this hazard analysis is to properly address the potential harm and risks
caused by the product to the user and the external environment.

\section{System Boundaries and Components}

\subsection{Hardware Components}
\subsubsection{Motion control modules}
The motion control module consists of a custom-designed PCB board utilizing an
STM32 microcontroller. This module is responsible for controlling the movement
of the camera gimbal. It receives commands from our Computer Vision module and
translates them into precise motor movements in real time.
\subsubsection{Camera gimbal}
The camera gimbal used in this project is a existing Camera Gimbal designed by
the McMaster Rocketry team. Since our team is not responsible for the design of
the gimbal, we can only estimate the potential hazards based on similar
products in the market.
\subsubsection{Jetson Orin Nano}
The Jetson Orin Nano is small, lightweight AI computer designed by NVIDIA. It
will be used to deploy our computer vision moodule and communicate directly
with the motion control module.
\subsubsection{Camera}
The camera is a off the shelf camera that is compatible with the Jetson Orin
Nano. It is able to provide 4k video input in real time.

\subsection{Software Components}
\subsubsection{Computer Vision Module}
Our computer vision module will be developed in Python. The base architecture
will be torch and pandas. It will be accerlerated with TensorRT in Jetson to
achieve real time stable performance.
\subsubsection{User Interface Module}
The user interface module will be developed using React. It will provide a
web-based interface for users to interact with the system, view camera feeds,
and control the camera gimbal.

\wss{Dividing the system into components will help you brainstorm the hazards.
  You shouldn't do a full design of the components, just get a feel for the major
  ones.  For projects that involve hardware, the components will typically include
  each individual piece of hardware.  If your software will have a database, or an
  important library, these are also potential components.}

\section{Critical Assumptions}

\subsection{Assumptions about the Operating Environment}
The system will not operating in any extreme environment (e.g., extreme
temperature, acidic environments, etc.). The system will be operating within 10
to 25 degree Celsius. The system will operate with a stable power supply.
\subsection{Assumptions about the User}
Anyone above 14 years old who uses electronic and exposed to software interface
on a daily basis for over 2 years with good intent.
\subsection{Assumptions about general of the shelf hardware components}
All off-the-shelf hardware components will function as advertised and will be
compatible with our system. This includes the camera, gimbal, Jetson Orin Nano
etc.
\subsection{Assumptions about the software libraries}
All software libraries used will not have any unknown security vulnerabilities.

\wss{These assumptions that are made about the software or system.  You should
  minimize the number of assumptions that remove potential hazards.  For instance,
  you could assume a part will never fail, but it is generally better to include
  this potential failure mode.}

\newpage
\section{Failure Mode and Effect Analysis}
\begin{table}[htbp]
  \caption{FMEA Table}\label{TblFMEALegend}
  \scalebox{0.6}{
    \begin{tblr}{
        colspec = {p{2cm} p{4cm} p{7cm} p{7cm} p{7cm}},
        row{1} = {font=\bfseries},
        hlines
      }
      Design Function                         & Failure Mode & Effect & Cause of Failure & Recommended Action \\
      Gimbal turn as CV control               &
      \begin{enumerate}
        \item Gimbal does not move
        \item Gimbal snaps
        \item Gimbal does not turn properly
      \end{enumerate}     &
      \begin{enumerate}
        \item \begin{enumerate}
                \item Camera loses the rocket
                \item user unable to debug the rocket launch
              \end{enumerate}
        \item \begin{enumerate}
                \item Camera loses the rocket
                \item Camera injures user (minor cuts and bruises)
                \item Camera and gimbal damages other property
                \item Camera, gimbal, motion control module, Jetson Orin Nano damages
              \end{enumerate}
        \item \begin{enumerate}
                \item Camera overshoots or loses the rocket
                \item Camera unable to track the rocket smoothly with high quality image
              \end{enumerate}
      \end{enumerate}
                                              &
      \begin{enumerate}
        \item \begin{enumerate}
                \item failed in communication between the motion control module and the gimbal
                \item failed in communication then CV module and the motion control module
                \item firmware bug in the motion control module or compatibility issue
                \item fried electric components in the motion control module or gimbal
              \end{enumerate}
        \item \begin{enumerate}
                \item improper instruction from the motion control module
                \item rocket angular velocity to the gimbal excesses the design limit
              \end{enumerate}
        \item \begin{enumerate}
                \item improper instruction from the motion control module
                \item CV module fails to provide proper instruction to the motion control module
                \item CV module lost the rocket
                \item damaged components in the motion control module or gimbal
              \end{enumerate}
      \end{enumerate}
                                              &
      \begin{enumerate}
        \item \begin{enumerate}
                \item careful handling in transportation and assembly to avoid damage to the
                      components
                \item Integration test before launch to ensure proper communication between all
                      modules
              \end{enumerate}
        \item \begin{enumerate}
                \item Set threshold of the maximum angular velocity to avoid gimbal snapping
                \item Set threshold of the maximum angle to avoid gimbal snapping
                \item Calculate the expected angular velocity of the camera system and the (SR1)
                      target to ensure they are within the design limit of the system
                \item safety goggles should be worn when operating the system (SR2)
              \end{enumerate}
        \item \begin{enumerate}
                \item design CV module to be robust in different environment
                \item design motion control module and CV module speed to have self-correcting (SR3)
                      mechanism
                \item design multiple tracking algorithms to avoid losing of the rocket

              \end{enumerate}
      \end{enumerate}
      \\
      User Interface control monitor tracking &
      \begin{enumerate}
        \item UI does not load
        \item UI loads but does not display camera feed
        \item User unable to understand and operate the UI
      \end{enumerate}
                                              &
      \begin{enumerate}
        \item \begin{enumerate}
                \item User unable to control the gimbal and camera
                \item Failed to monitor the rocket launch
                \item Failed to track the rocket launch
                \item User frustration
              \end{enumerate}
        \item \begin{enumerate}
                \item User unable to monitor rocket launch in real time
                \item User unable to track the rocket launch in real time
                \item User unable to select desired tracking mode and target
              \end{enumerate}
        \item \begin{enumerate}
                \item User unable to control the gimbal and camera
                \item Failed to monitor the rocket launch
                \item Failed to track the rocket launch
                \item User frustration
              \end{enumerate}
      \end{enumerate}
                                              &
      \begin{enumerate}
        \item \begin{enumerate}
                \item Bug in the UI software
                \item Issue in CD pipeline
                \item Browser compatibility issue
                \item OS environment issue
              \end{enumerate}
        \item \begin{enumerate}
                \item Issue in API between UI and CV module
                \item Camera or Jetson Orin Nano hardware issue
              \end{enumerate}
        \item \begin{enumerate}
                \item Poorly designed UI
                \item Lack of user manual and documentation
              \end{enumerate}
      \end{enumerate}
                                              &
      \begin{enumerate}
        \item \begin{enumerate}
                \item Thorough testing in different browsers and OS environment
                \item Proper CI/CD pipeline to ensure smooth deployment
                \item Provide user manual and documentation for installation and troubleshooting
              \end{enumerate}
        \item \begin{enumerate}
                \item Thorough testing of the API between UI and CV module
                \item Provide user manual and documentation for installation and troubleshooting
                \item Careful handling in transportation and assembly to avoid damage to the
                      components
              \end{enumerate}
        \item \begin{enumerate}
                \item User testing and feedback to improve UI design
                \item Provide user manual and documentation for installation and troubleshooting
                \item Provide user with access to an AI Agent that read the user manual and
                      documentation
              \end{enumerate}
      \end{enumerate}

    \end{tblr}
  }
\end{table}

\wss{Include your FMEA table here. This is the most important part of this document.}
\wss{The safety requirements in the table do not have to have the prefix SR.
  The most important thing is to show traceability to your SRS. You might trace to
  requirements you have already written, or you might need to add new
  requirements.}
\wss{If no safety requirement can be devised, other mitigation strategies can be
  entered in the table, including strategies involving providing additional
  documentation, and/or test cases.}

\newpage
\section{Safety and Security Requirements}

\begin{enumerate}
  \item SR1: The system will be placed in a distance which the expected angular
        velocity will not exceed the design limit of the gimbal.
  \item SR2: all operators must wear safety goggles when operating the system.
  \item SR3: The system will have a self-correcting mechanism to ensure the camera can
        track the rocket smoothly in real time.
\end{enumerate}

\section{Roadmap}

\subsection{Capstone Implementation Plan (Sprints \& Milestones)}
\begin{table}[H]
  \caption{Capstone Roadmap (8 sprints)}\label{tab:capstone-roadmap}
  \renewcommand{\arraystretch}{1.15}
  \begin{tabularx}{\textwidth}{C{1.2cm} L{4.2cm} Y Y Y}
    \toprule
    \textbf{Sprint}                                                                                                               & \textbf{Primary Goals} & \textbf{Deliverables} & \textbf{Safety Focus (Traceability)} & \textbf{Verification / Evidence} \\
    \midrule
    1                                                                                                                             &
    System bring-up: Jetson Orin Nano, camera, repo CI. Define interfaces (UI \(\leftrightarrow\) CV \(\leftrightarrow\) Motion). &
    Env setup scripts; interface spec; smoke tests.                                                                               &
    Baseline hazard review; confirm assumptions (power, temp).                                                                    &
    Checklist, boot logs, interface mock tests.                                                                                                                                                                                                              \\
    \midrule
    2                                                                                                                             &
    Motion control MVP: gimbal command path; encoder/limits readout.                                                              &
    Motor driver config; limit-switch handling; teleop tool.                                                                      &
    SR1 (angular velocity/angle limits) draft.                                                                                    &
    Bench test: limit cut-offs; plotted step responses.                                                                                                                                                                                                      \\
    \midrule
    3                                                                                                                             &
    CV tracking MVP on sample footage (offline \(\rightarrow\) near-real-time).                                                   &
    Tracker module; target reacquisition logic.                                                                                   &
    SR3 draft (self-correction, loss recovery).                                                                                   &
    Replay tests; failure-injection (occlusion) logs.                                                                                                                                                                                                        \\
    \midrule
    4                                                                                                                             &
    UI prototype: live feed, mode select, status \& alarms.                                                                       &
    Web UI with HUD; error banners; safe-state button.                                                                            &
    FMEA controls for “UI not loading/displaying”.                                                                                &
    Cross-browser tests; watchdog for API health.                                                                                                                                                                                                            \\
    \midrule
    5                                                                                                                             &
    Closed-loop integration: CV \(\rightarrow\) Motion with rate limiting.                                                        &
    End-to-end demo; rate limiter \& soft stops.                                                                                  &
    SR1 formalized (hard/soft limits); map FMEA “gimbal snap”.                                                                    &
    Angular-rate sweep; saturation tests; limit breach proofs.                                                                                                                                                                                               \\
    \midrule
    6                                                                                                                             &
    Robustness: fallback trackers, loss detection, auto-recenter.                                                                 &
    Multi-strategy tracker; loss timers; homing routine.                                                                          &
    SR3 formalized (self-correction).                                                                                             &
    Scenario matrix (lighting/occlusion); recovery time stats.                                                                                                                                                                                               \\
    \midrule
    7                                                                                                                             &
    Safety hardening: interlocks, E-stop, checklists, SOP.                                                                        &
    E-stop wiring/command; pre-launch checklist; SOP v1.                                                                          &
    SR2 (PPE), SR1 (interlocks), UI hazard cues.                                                                                  &
    Dry-run with operators; checklist sign-offs; video evidence.                                                                                                                                                                                             \\
    \midrule
    8                                                                                                                             &
    Verification \& docs: finalize FMEA, traceability, test report.                                                               &
    FMEA v2; SRS trace matrix; safety test report; demo video.                                                                    &
    All SRs: SR1–SR3 complete and traced.                                                                                         &
    Pass/fail summary; residual-risk log; release tag.                                                                                                                                                                                                       \\
    \bottomrule
  \end{tabularx}
\end{table}

\subsection{Milestones}
\begin{enumerate}
  \item M1 (End Sprint 2): Motion MVP with enforceable soft limits (SR1 draft).
  \item M2 (End Sprint 4): UI prototype with health/status and safe-state control.
  \item M3 (End Sprint 6): Closed-loop tracking with recovery/self-correction (SR3).
  \item M4 (End Sprint 8): Safety package complete (FMEA v2, SR1–SR3 verified, docs).
\end{enumerate}

\subsection{Traceability Plan}
\begin{itemize}
  \item FMEA \(\rightarrow\) SR mapping: “Gimbal snaps/overshoots” \(\rightarrow\) SR1;
        “Tracking lost” \(\rightarrow\) SR3; “UI unavailable/unclear” \(\rightarrow\)
        UI alarms, SOP, tests.
  \item Each sprint closes with: (a) updated FMEA row(s), (b) test artifacts linked to
        SR IDs, (c) residual risk review.
\end{itemize}

\subsection{Acceptance Criteria}
\begin{itemize}
  \item \textbf{SR1}: Max angular velocity/angle never exceeded in any commanded or fault path; evidence: logs + plotted limits across stress tests.
  \item \textbf{SR2}: Operators use PPE; SOP and checklist enforced; evidence: signed checklist + demo footage.
  \item \textbf{SR3}: From target-loss events, system recovers or safely recenters within \(T_{\max}\) seconds in \(\geq 95\%\) of trials.
\end{itemize}

\subsection{Post-Capstone Backlog (Future Work)}
\begin{enumerate}
  \item Hardware redundancy (limit sensors, power rails), formal E-stop certification.
  \item Model optimization (TensorRT INT8), thermal profiling, enclosure/EMI review.
  \item Expanded UI: guided wizards, operator training mode, automated pre-flight.
  \item Data logging for incident analysis; automated regression test farm.
  \item Security hardening: signed releases, SBOM, dependency scanning.
\end{enumerate}
\newpage{}

\section*{Appendix --- Reflection}

\wss{Not required for CAS 741}

The purpose of reflection questions is to give you a chance to assess your own
learning and that of your group as a whole, and to find ways to improve in the
future. Reflection is an important part of the learning process.  Reflection is
also an essential component of a successful software development process.  

Reflections are most interesting and useful when they're honest, even if the
stories they tell are imperfect. You will be marked based on your depth of
thought and analysis, and not based on the content of the reflections
themselves. Thus, for full marks we encourage you to answer openly and honestly
and to avoid simply writing ``what you think the evaluator wants to hear.''

Please answer the following questions.  Some questions can be answered on the
team level, but where appropriate, each team member should write their own
response:


\begin{enumerate}
  \item What went well while writing this deliverable?
  \item What pain points did you experience during this deliverable, and how did you
        resolve them?
  \item Which of your listed risks had your team thought of before this deliverable,
        and which did you think of while doing this deliverable? For the latter ones
        (ones you thought of while doing the Hazard Analysis), how did they come about?
  \item Other than the risk of physical harm (some projects may not have any
        appreciable risks of this form), list at least 2 other types of risk in
        software products. Why are they important to consider?
\end{enumerate}

\begin{enumerate}
  \item \textbf{What went well while writing this deliverable? (CV perspective)} \
        From the computer vision side, what went well was how the hazard analysis helped us think beyond just performance and accuracy. Normally, we focus on improving tracking or inference speed, but this deliverable made us consider how our module’s failures could affect safety and the physical system. For example, if the CV model misidentifies or loses the rocket, the gimbal might swing too far or too fast. Listing these in the FMEA table showed clear links between the vision system and hardware movement. It also encouraged us to add fallback ideas, like recovery tracking or self-correction, which became part of SR3. The process improved our system awareness — we now design the CV pipeline not only to be fast, but also predictable and safe.

  \item \textbf{What pain points did you experience during this deliverable, and how did you resolve them? (Development/UI perspective)} \
        From the development side, one major challenge was figuring out how to document software-related hazards in a meaningful way. It was easier to picture a gimbal breaking than a user interface “failing,” so at first, the risks in our module sounded less serious. We realized, however, that a software crash or unclear UI could lead to real operational problems if users can’t control or monitor the camera. Formatting the FMEA table in LaTeX was also tricky — the entries for causes and effects got very long. We solved this by using short numbered lists and grouping related issues under one failure mode. Collaborating with the hardware and CV members also helped, because they could explain how our software might contribute to physical or timing failures, which gave us better context for our part of the analysis.

  \item \textbf{Which of your listed risks had your team thought of before this deliverable, and which did you think of while doing this deliverable? How did they come about? (Hardware perspective)} \
        From the hardware perspective, we had already identified the obvious mechanical hazards before writing this document — such as motor overdriving, gimbal snapping, or loose mounts. These were discussed early in the semester when testing servo drivers. However, new risks came up during the hazard analysis that we hadn’t thought about before. For instance, we realized that communication loss between the motion control board and the Jetson could cause unpredictable movements, and that power fluctuations might damage sensitive components. These ideas surfaced while building the FMEA table, since it forced us to trace every “what if” scenario from signal input to physical motion. It also made us think about how environmental assumptions (like stable temperature and power) affect reliability. Writing this section helped the hardware team plan safer wiring, limit switches, and emergency stops for future testing.

  \item \textbf{Other than the risk of physical harm, list at least 2 other types of risk in software products. Why are they important to consider? (QA perspective)} \
        From a quality assurance perspective, two major risks beyond physical harm are \textit{data integrity risk} and \textit{maintainability risk}. Data integrity risk means that logs or sensor readings might be lost, corrupted, or inconsistent — which can make debugging almost impossible. If the system fails but there’s no reliable data record, identifying the cause becomes guesswork. Maintainability risk comes from unclear code structure, missing documentation, or complex dependencies that make it hard to update safely. These are important to consider because a project that cannot be maintained or verified quickly becomes unsafe, even if it appears to “work.” Good QA practices like regression testing, version control, and traceable requirements help prevent these issues. Recognizing these risks early will help us build a more dependable system over time.
\end{enumerate}

\end{document}